\subsection{Discrete valuations}

DEFINITION 3. Let K be a (not necessarily commutative) field, v a valuation on K and \(\Gamma\) the order group of v. v is called discrete if there exists a (necessarily unique) isomorphism of the ordered group \(\Gamma\) onto Z. Let \(\gamma\) be the element of \(\Gamma\) corresponding to 1 under this isomorphism; every element u of K such that \(v(u) = \gamma\) is called a uniformizer of v. A discrete valuation is called normed if its order group is Z.

For example the valuation \(v_{p}\) defined by an extremal element p of a principal ideal *or factorial* domain is a normed discrete valuation which admits p as a uniformizer. In particular, if k is a field, \(k[[T]]\) is the ring of a discrete valuation on \(k((T))\) which admits T as a uniformizer. * Let S be a connected complex analytic variety of dimension 1, K the field of meromorphic functions on S and \(z_{0}\) a point of S; the set of \(f \in K\) which are holomorphic at \(z_{0}\) is the ring of a discrete valuation v; for a function \(f \in K\) to be uniformizing for v, it is necessary and sufficient that it be holomorphic and zero at \(z_{0}\) and that there exist a neighbourhood V of \(z_{0}\) in S such that the restriction of f to V be a homomorphism of V onto a neighbourhood of the origin in C. It is this example and other analogues which are the origin of the word ``uniformizer''.*

PROPOSITION 8. Let K be a (not necessarily commutative) field, v a discrete valuation on K, A the ring of v and u a uniformizer for v. The non-zero ideals of A are two-sided and of the form Au \(^{n}\) (n ≥ 0).

It may be assumed that \(v\) is normed, so that \(v(u) = 1\) . For all \(x \in \mathbf{K}^*\) , there is an integer \(n \in \mathbf{Z}\) such that \((v(x) = n = v(u^n))\) and hence we may write

\[
x = z u ^ {n} = u ^ {n} z ^ {\prime},
\]

where \(z, z'\) are two invertible elements of the ring A; whence the proposition.

PROPOSITION 9. Let A be a local integral domain distinct from its field of fractions. The following conditions are equivalent:

\begin{enumerate}
\def\labelenumi{(\alph{enumi})}
\item
  A is the ring of a discrete valuation.
\item
  A is a principal ideal domain.
\item
  The ideal \(\mathfrak{m}(\mathbf{A})\) is principal and \(\bigcap_{n=1}^{\infty} \mathfrak{m}(\mathbf{A})^n = (0)\) .
\item
  A is a Noetherian ring and \(\mathfrak{m}(\mathbf{A})\) is principal.
\item
  A is a Noetherian valuation ring.
\end{enumerate}

Proposition 8 shows that (a) implies (b), (d) and (e). If A is a principal ideal domain, then \(\mathfrak{m}(\mathbf{A}) = \mathbf{A}\mathbf{u}\) and every non-zero ideal of A is of the form \(\mathbf{A}\mathbf{u}^n\) since A is local (Algebra, Chapter VII, §1, no. 3, Theorem 2); therefore \(\bigcap_{n=1}^{\infty} \mathfrak{m}(A)^n = 0\) ; this shows that (b) implies (c). On the other hand (d) implies (c) (Chapter III, §3, no. 2, Corollary to Proposition 5); by Proposition 2 of §1, no. 4, (c) implies (a). Thus conditions (a), (b), (c), (d) are equivalent and imply (e). Finally suppose (e) holds and let us show that (b) holds; it will be sufficient to prove the following lemma:

LEMMA 1. Let A be a valuation ring. Every finitely generated torsion-free A-module is free. Every finitely generated ideal of A is principal. Every torsion-free A-module is flat.

Let E be a finitely generated torsion-free A-module and let \(x_{1}, \ldots, x_{n}\) be generators of E which are minimal in number; we show that they are linearly independent. If \(\sum_{i=1}^{n} a_i x_i = 0\) ( \(a_i \in A\) ) is a non-trivial relation between the \(x_i\) , one of the \(a_i\) , say \(a_1\) , divides all the others since the set of principal ideals of \(A\) is totally ordered by inclusion (§ 1, no. 2, Theorem 1); then \(a_1 \neq 0\) since the relation is non-trivial. As \(E\) is torsion-free, we can divide by \(a_1\) , which amounts to assuming that \(a_1 = 1\) . But then \(x_1\) is a linear combination of \(x_2, \ldots, x_n\) , contrary to the minimal character of \(n\) . Hence \(E\) is free.

In particular every finitely generated ideal a of A is principal, all the elements of a system of generators of a being multiples of one of them. Proposition 3 of Chapter I, § 2, no. 4 then shows that every torsion-free A-module is flat.

